\paragraph{Propagating attributions through a pipeline.} Redistributing an attribution from a stage's outputs to its inputs, in proportion to each input's contribution, is an established idea. DeepLIFT's chain rule for multipliers~\cite{shrikumar2017learning} and layer-wise relevance propagation~\cite{bach2015pixel} do this layer by layer inside a network, and G-DeepSHAP~\cite{chen2022explaining} lifts it to a series of models: its generalized rescale rule propagates Shapley values from one model to the next, satisfies efficiency at every stage, and is exact for chains of linear models. Our pullback $b = Pa$ coincides with that rule when $P$ holds each byte's share of each feature's change, and we claim no novelty for the propagation step itself. What these methods need is that every stage be a model with its own efficiency-satisfying attribution method (a network, a tree ensemble, a linear map). A real feature extractor is a program: a parser, hashing, string scans, windowed entropy, and header lookups over discrete, variable-length input with no natural baseline. We construct the stage-wise matrix for such a program by instrumenting it, estimate it without instrumentation when that is not possible, and characterize what any transfer through it can preserve from the matrix alone. In the other direction, virtual inspection layers~\cite{vielhaben2024explainable} move relevance through an \emph{invertible} linear transform (time to frequency) into a domain the model never computes; the pushforward does this through a non-invertible program, into the feature vocabulary of a detector that never reads features. Aggregating attributions over pixels into segments or tokens into words~\cite{jullum2021groupshapley} is the special case of a 0/1 provenance within one space, which we use when coarsening $P$ to windows of bytes.

\paragraph{Explaining malware detectors.} Explanations of malware detectors follow the two representations. Detectors over hand-crafted features are explained at the feature level~\cite{fan2020can, wu2021android, zhu2024dynamic, warnecke2020evaluating}, most often with SHAP over EMBER features; Mohseni et al.~\cite{mohseni2026shap} argue that such explanations are unreliable on their own, because interventional SHAP queries off-manifold coalitions and redundant features split credit. Detectors over raw input are explained over the input itself: activations and gradients show that MalConv relies heavily on the header~\cite{coull2019activation, bose2020explaining, demetrio2019explaining}, and byte-level SHAP on MalConv, aggregated over PE regions, has been used to choose where to inject adversarial content~\cite{aryal2024explainability}. FINER~\cite{he2023finer} formalizes the path from problem space to feature space to explanation for security classifiers and explicitly leaves out pattern-driven detectors such as EMBER-based ones, arguing that hashing and aggregation lose too much information to map their explanations back to program elements. Our work addresses exactly this case: $P$ maps an EMBER detector's explanations onto the bytes, and its statistics measure, rather than assume, how much of an explanation can be placed in the file. Problem-space explanations can also be checked with real files, which is how we evaluate the pullback (Section~\ref{sec:XML:Pullback}); this complements, but does not resolve, the feature-level concerns of Mohseni et al.

\paragraph{From feature explanations to the binary.} Attacks on malware detectors must turn a feature-space change into a valid file, the inverse feature-mapping problem of Pierazzi et al.~\cite{pierazzi2020intriguing}. Explanation-guided attacks solve it by hand: Severi et al.~\cite{severi2021explanation} select backdoor triggers from SHAP values on EMBER features and realize them with a purpose-built PE editing utility, keeping only the 17 features they could map to file content; Rosenberg et al.~\cite{rosenberg2020generating} modify features one at a time with type-specific transformations; and toolkits of functionality-preserving manipulations~\cite{demetrio2021adversarial} realize changes at known PE locations. Mohammed and Zhang~\cite{mohammed2026beyond} find that a static classifier's top features are packing and metadata artifacts and confirm it by manually analyzing the bytes of PE files. In each case, a person decides which bytes a feature comes from. $P$ records this mapping for every feature and every file as a by-product of extraction; we use it defensively, to compare and audit detectors, not to build attacks.

\paragraph{Detectors with different feature sets.} Production antivirus engines combine several models over different representations of the same file, including models over featurized metadata and deep networks over full file contents~\cite{microsoft2019insideout}. The 2019 machine-learning evasion competition targeted exactly our pair, MalConv and a LightGBM model on EMBER features~\cite{endgame2019mlsec}; the winning defense of its 2020 edition combined models over distinct feature sets so that an attacker must evade all of them at once~\cite{quiring2020against}; and adversarial binaries crafted for one detector often transfer poorly to another~\cite{salman2024robustness}. These works compare such detectors by their outputs. The pullback compares them by their evidence, per file and in one space: which bytes both rely on, and, when they disagree, whether the flagging detector reads content the other ignores or reads the other way (Section~\ref{sec:XML:Comparison}).

\paragraph{Shortcuts and planted ground truth.} Spurious correlations are a known pitfall of learning-based security systems~\cite{arp2022and}, and finding the ones a raw-byte detector rests on currently means reading byte-level saliency file by file~\cite{demetrio2019explaining, bose2020explaining} or explaining a surrogate~\cite{mohammed2026beyond}. The pushforward aggregates a raw-byte detector's per-file explanations into named, recurring reliance across a corpus, which we compare with a surrogate and confirm by scrambling the bytes $P$ implicates (Section~\ref{sec:XML:Audit}). Planted triggers are an established ground truth for evaluating explainers~\cite{lin2021what}, alongside deletion- and retraining-based metrics~\cite{petsiuk2018rise, hooker2019benchmark} and sanity checks~\cite{adebayo2018sanity}; we plant triggers in the problem space of real PE files, so that the ground truth is a set of bytes.

\paragraph{Provenance.} In databases, provenance records which input tuples produce an output and how, with a compositional algebra~\cite{green2007provenance}; dynamic taint tracking records the same at the level of program instructions. Our provenance is of this kind: a weighted record of which input bytes produce each feature, obtained here by instrumenting the extractor (taint tracking is one way to automate it). It is unrelated to the system provenance graphs used for intrusion detection~\cite{cheng2024kairos, milajerdi2019poirot}, which record causal relationships among processes, files, and sockets on a host.
